\exercisegroup

\begin{enumerate}
\def\labelenumi{\arabic{enumi})}
\item
  Let X be a Lindelöf space (TG, IX, appendix 1, p. 75, def. 1) that is loopable. Show that, for every point \(x \in X\) , the Poincaré group \(\pi_1(X, x)\) is countable.
\item
  Let Y be a uniform space and let X be a subspace of Y. Let \((\mathrm{U}_{i})_{i\in\mathrm{I}}\) be an open cover of X.
\end{enumerate}

\begin{enumerate}
\def\labelenumi{\alph{enumi})}
\tightlist
\item
  Prove that there exists an open neighborhood V of X in Y and an open covering \((V_{i})_{i\in I}\) of V satisfying the following properties:
\end{enumerate}

\begin{enumerate}
\def\labelenumi{(\roman{enumi})}
\item
  For every \(i\) , \(\mathrm{U}_i = \mathrm{V}_i \cap \mathrm{X}\) ;
\item
  For every subset J of I such that \(\bigcap_{j\in J}V_j\neq \varnothing\), we have \(\bigcap_{j\in J}U_j\neq \varnothing .\)
\end{enumerate}

\begin{enumerate}
\def\labelenumi{\alph{enumi})}
\setcounter{enumi}{1}
\item
  Assume that Y is locally connected and that the \(U_{i}\) are connected. Prove that there exists such a covering in which the \(V_{i}\) are connected open subsets of Y.
\item
  Suppose that the \(U_{i}\) are connected and let, for every \(i \in I\), a point \(x_{i} \in U_{i}\). Suppose that the image of \(\pi_{1}(U_{i}, x_{i})\) in \(\pi_{1}(X, x_{i})\) is reduced to the identity element. Show that there exists a neighborhood V of X in Y such that, for every \(x \in X\), the canonical homomorphism \(\pi_{1}(X, x) \to \pi_{1}(V, x)\) admits a retraction.
\end{enumerate}

\begin{enumerate}
\def\labelenumi{\arabic{enumi})}
\setcounter{enumi}{2}
\tightlist
\item
  \begin{enumerate}
  \def\labelenumii{(\arabic{enumii})}
  \setcounter{enumii}{2}
  \tightlist
  \item
    Let \(X\) be a topological space and let \(x\) be a point of \(X\) .
  \end{enumerate}
\end{enumerate}

Let G be a group and let \(\varphi\colon\pi_{1}(\mathrm{X},x)\to\mathrm{G}\) be a group homomorphism. Let \((\mathrm{U}_{n})_{n\in\mathbf{N}}\) be a decreasing sequence of subspaces of X whose intersection is reduced to the point \(x\) . For all \(n\) , let \(\mathrm{G}_{n}\) be the set of elements of G that are images under \(\varphi\) of the class of a loop contained in \(\mathrm{U}_{n}\) .

\begin{enumerate}
\def\labelenumi{\alph{enumi})}
\item
  If G is countable, the sequence \((\mathrm{G}_{n})\) is stationary. (Reduce to the case where \((\mathrm{G}_{n}:\mathrm{G}_{n+1})\geqslant n+1\) for all n and consider an enumeration \((\gamma_{n})_{n}\) of G; construct recursively elements \(\lambda_{n}\in G_{n}\) such that, on writing \(w_{n}=\lambda_{1}\ldots\lambda_{n}\), one has \(\gamma_{i}\notin w_{n}G_{n+1}\) for \(1\leqslant i\leqslant n\). Show then that the intersection \(\bigcap_{i}w_{i}G_{i+1}\) is not empty.)
\item
  Suppose that G is abelian and that its identity element is the only infinitely divisible element of G. Then \(\bigcap_{n} G_{n}\) is reduced to the identity element. (Let a be an element of the intersection; for every n, let \(\gamma_{n} \in \pi_{1}(U_{n}, x)\) whose image under \(\varphi\) is equal to a. Construct by induction path classes \(\lambda_{n} \in \pi_{1}(U_{n}, x)\) such that \(\lambda_{n} = \gamma_{n}\lambda_{n+1}^{d}\), where d is a fixed integer \(\geqslant 2\). Show that \(b = (d - 1)\varphi(\lambda_{1}) + a\) is infinitely divisible. Deduce that a = 0.)
\item
  More generally, if G is abelian and every infinitely divisible element of G is a torsion element, then the intersection of the \(G_{n}\) consists of torsion elements.
\end{enumerate}

(3) The results of this exercise and the following one are taken from the article by J. W. CANNON and G. R. CONNER, “On the fundamental groups of one-dimensional spaces,” Topology and its Applications 153 (2006), pp. 2648–2672.

\begin{enumerate}
\def\labelenumi{\arabic{enumi})}
\setcounter{enumi}{3}
\tightlist
\item
  Let X be a connected and locally connected by arcs topological space. Let x be a point of X that has a countable fundamental system of neighborhoods. Let S be a set and let \(\varphi\colon\pi_{1}(\mathrm{X},x)\to\mathbf{Z}^{(\mathrm{S})}\) be a surjective group homomorphism.
\end{enumerate}

\begin{enumerate}
\def\labelenumi{\alph{enumi})}
\item
  Prove that S is countable.
\item
  Assume that X is compact; prove that S is finite.
\item
  We assume that \(\pi_{1}(X,x)\) is a free abelian group; prove that X is loopable.
\item
  We assume that \(\pi_{1}(X,x)\) is a free group; prove that X is loopable. (Use the fact that the intersection of the subgroups of the lower central series of a free group is reduced to the identity element.)
\end{enumerate}

\begin{enumerate}
\def\labelenumi{\arabic{enumi})}
\setcounter{enumi}{4}
\tightlist
\item
  Let P be the Hawaiian earring (III, p. 336, exerc. 6, whose notation we adopt). For \(n \in \mathbf{N}^*\) , let \(\alpha_n \in \pi_1(\mathrm{P}, o)\) be the class of a loop in \(\mathrm{C}_n\) whose class generates \(\pi_1(\mathrm{C}_n, 0)\) .
\end{enumerate}

\begin{enumerate}
\def\labelenumi{\alph{enumi})}
\item
  Prove that the map \(\mathrm{Hom}(\pi_{1}(\mathrm{P},0),\mathbf{Z})\to\mathbf{Z}^{\mathbf{N}}\) , \(\varphi\mapsto(\varphi(\alpha_{n}))_{n\in\mathbf{N}}\) is an injective group homomorphism and that its image is equal to \(\mathbf{Z}^{(\mathbf{N})}\) . (4)
\item
  The group \(\pi_{1}(P,0)\) is not a free group; the quotient of \(\pi_{1}(P,0)\) by its derived group is not a free abelian group.
\item
  Let A be the Hawaiian archipelago (III, p. 338, exerc. 10). Prove that every homomorphism from \(\pi_1(A, o)\) into \(\mathbf{Z}\) is trivial.
\end{enumerate}

(4) See B. DE SMIT, “The fundamental group of the Hawaiian earring is not free,” International Journal of Algebra and Computation, Vol. 2, No. 1 (1992), pp. 33–37.

